\chapter{ANALYSIS/DISCUSSION}
\label{chapter5}
This chapter presents the interpretation and discussion of the results obtained in the previous chapter. The purpose of this chapter is to analyze the findings, explain observed trends, and relate them to the research objectives and existing literature.

This chapter is organized according to the main objectives of the study. Each section discusses the corresponding results and provides interpretation based on theoretical concepts and previous research.

% -------------------------------------------------
\section{General Discussion}

This section provides an overall interpretation of the results presented in Chapter 4. The findings are analyzed in relation to the research objectives outlined in Chapter 1.

Rather than simply restating the results, this section explains the significance of the findings. The relationships between variables, observed patterns, and deviations from expected behavior should be discussed.

% -------------------------------------------------
\section{Discussion of Topic 1}

This section discusses the results related to Topic 1.

\subsection{Case 1 Analysis}

The results from Case 1 are analyzed in detail. Explain why the observed outcomes occurred and whether they align with theoretical expectations or previous studies.

\subsection{Case 2 Analysis}

Compare the results obtained in Case 2 with those of Case 1. Highlight improvements, differences, or unexpected observations.

For example:
\begin{itemize}
    \item Explain variations in performance
    \item Identify contributing factors
    \item Relate findings to methodology
\end{itemize}

% -------------------------------------------------
\section{Discussion of Topic 2}

This section presents analysis related to Topic 2. Interpret the results by comparing them with similar studies discussed in Chapter 2.

Explain whether the results support or contradict existing literature. If differences are observed, provide possible reasons such as methodology differences, assumptions, or limitations.

% -------------------------------------------------
\section{Discussion of Topic 3}

This section focuses on the analysis of results corresponding to Topic 3. Discuss significant findings and explain their implications.

Highlight any new insights gained from the study and explain how they contribute to the research area.

% -------------------------------------------------
\section{Comparison with Existing Literature}

This section compares the findings of this study with results from previous research.

\begin{table}[!tbh]
\centering
\caption{Comparison of current results with existing literature}
\label{tab:comparison_literature}
\begin{tabular}{|c|c|c|}
\hline
\textbf{Study} & \textbf{Result} & \textbf{Remarks} \\
\hline
Previous Study A & Result A & Limitation observed \\
Previous Study B & Result B & Partial agreement \\
Proposed Study & Result C & Improved performance \\
\hline
\end{tabular}
\vspace{0.3em}
\end{table}

Discuss how the current study differs from or improves upon existing works. Clearly indicate the contribution of the research.

% -------------------------------------------------
\section{Implications of the Study}

This section highlights the practical and theoretical implications of the findings.

\begin{itemize}
    \item Practical significance of the results
    \item Contribution to the research field
    \item Potential applications of the findings
\end{itemize}

Explain how the results can be used in real-world scenarios or future research.

% -------------------------------------------------
\section{Limitations of the Study}

This section discusses limitations encountered during the study. These may include:

\begin{itemize}
    \item Constraints in data or resources
    \item Limitations of methodology
    \item Assumptions affecting results
\end{itemize}

Acknowledging limitations helps in evaluating the reliability and applicability of the research findings.